\documentclass[12pt,a4paper]{amsart}
\usepackage{amsmath,amssymb,amsthm}
\usepackage{mathtools}
\usepackage[utf8]{inputenc}
\usepackage[T1]{fontenc}
\usepackage{hyperref}
\usepackage{enumitem,booktabs,array}
\usepackage{geometry}
\geometry{margin=2.5cm}

% -------------------------------------------------------
% Theorem environments
% -------------------------------------------------------
\newtheorem{theorem}{Theorem}[section]
\newtheorem{lemma}[theorem]{Lemma}
\newtheorem{corollary}[theorem]{Corollary}
\newtheorem{proposition}[theorem]{Proposition}
\theoremstyle{definition}
\newtheorem{definition}[theorem]{Definition}
\newtheorem{example}[theorem]{Example}
\newtheorem{remark}[theorem]{Remark}
\newtheorem{conjecture}[theorem]{Conjecture}

% -------------------------------------------------------
% Custom commands
% -------------------------------------------------------
\newcommand{\N}{\mathbb{N}}
\newcommand{\Z}{\mathbb{Z}}
\newcommand{\Q}{\mathbb{Q}}
\newcommand{\R}{\mathbb{R}}
\newcommand{\C}{\mathbb{C}}
\DeclareMathOperator{\lcm}{lcm}
\DeclareMathOperator{\ord}{ord}

% -------------------------------------------------------
\begin{document}
% -------------------------------------------------------

\title{The Erd\H{o}s--Straus Conjecture:\\
       Egyptian Fractions, Diophantine Equations,\\
       and Partial Results}

\author{Michael Fuhrmann}
\address{Specialist Mathematics Project, Build 133}
\email{specialist-maths@localhost}
\date{\today}

\begin{abstract}
The Erd\H{o}s--Straus conjecture, proposed by Paul Erd\H{o}s and Ernst
Gabor Straus in 1948, asserts that for every integer $n \geq 2$ there exist
positive integers $a, b, c$ such that
\[
  \frac{4}{n} = \frac{1}{a} + \frac{1}{b} + \frac{1}{c}.
\]
This is a statement about Egyptian fractions (representations of rationals
as sums of unit fractions) and one of the simplest open problems in
elementary number theory.  The conjecture has been verified computationally
for all $n \leq 10^{14}$ and is known to hold for almost all~$n$ in the
density sense (Schinzel, 1956; Vaughan, 1970).  We present the historical
background, all major constructive approaches, modular reductions, density
results, the connection to Diophantine equations of the form
$4abc = n(ab+ac+bc)$, and a summary of computational verification for
all primes up to $10^9$.  The conjecture remains open.
\end{abstract}

\maketitle
\tableofcontents

% =========================================================
\section{Introduction}
% =========================================================

One of the oldest traditions in number theory is the representation of
rational numbers as sums of unit fractions $1/k$, $k \in \N$.  Such
representations were used by the ancient Egyptians and appear, for example,
in the Rhind Mathematical Papyrus (circa 1650~BCE).  A \emph{unit fraction}
or \emph{Egyptian fraction} is simply a fraction of the form $1/k$ with
$k$ a positive integer.

In 1948, Paul Erd\H{o}s and Ernst Gabor Straus posed the following question:
Is it true that every fraction $4/n$ can be written as a sum of exactly
three unit fractions?  This deceptively simple question has remained
unresolved for more than seventy years.

\begin{conjecture}[Erd\H{o}s--Straus, 1948]\label{conj:ES}
For every integer $n \geq 2$ there exist positive integers $a, b, c$ such that
\[
  \frac{4}{n} = \frac{1}{a} + \frac{1}{b} + \frac{1}{c}.
\]
\end{conjecture}

The difficulty lies in the universality of the claim: we need a
representation for \emph{every} $n$, not just for a large class of them.
While the conjecture has been checked computationally up to $n = 10^{14}$
without a single exception, a proof remains elusive.

The problem is emblematic of the genre of Diophantine problems that are
easy to state, easy to verify for individual values, yet extraordinarily
hard to settle in full generality.  It sits alongside the Goldbach
conjecture and the Collatz conjecture as one of the most famous unsolved
problems accessible to non-specialists.

\subsection*{Organisation of this paper}

Section~2 reviews the history of Egyptian fractions.  Section~3 gives the
formal statement and equivalent reformulations.  Section~4 develops the
three main constructive approaches.  Section~5 treats modular arithmetic
reductions.  Section~6 covers density and analytic results.
Section~7 summarises computational verification.  Section~8 lists
related open problems.  Section~9 contains the references.

% =========================================================
\section{Egyptian Fractions: History and Foundations}
% =========================================================

\subsection{Ancient origins}

The Egyptians did not write fractions $p/q$ with $p > 1$; instead they
expressed every rational number as a finite sum of \emph{distinct} unit
fractions.  The Rhind Papyrus (c.\ 1650~BCE) contains a table of
decompositions of $2/n$ for odd $n$ from $3$ to $101$.  For example:
\[
  \frac{2}{5} = \frac{1}{3} + \frac{1}{15}, \quad
  \frac{2}{7} = \frac{1}{4} + \frac{1}{28}, \quad
  \frac{2}{13} = \frac{1}{8} + \frac{1}{52} + \frac{1}{104}.
\]
This tradition was carried on by Fibonacci (Leonardo of Pisa), who in his
\emph{Liber Abaci} (1202) described a greedy algorithm --- now called the
\emph{Fibonacci--Sylvester algorithm} --- for expanding any rational
number into unit fractions.

\subsection{The Fibonacci--Sylvester greedy algorithm}

\begin{theorem}[Fibonacci--Sylvester algorithm]\label{thm:sylvester}
Every positive rational number $r = p/q$ with $0 < p/q \leq 1$ has a
finite representation as a sum of distinct unit fractions.
\end{theorem}

The algorithm works by iterating: at each step replace $p/q$ by
$p/q - 1/\lceil q/p \rceil$.  The crucial point is that this process
terminates, since the numerator strictly decreases at each step.

\begin{example}
We expand $4/7$ using the greedy algorithm.
\[
  \frac{4}{7}: \quad \left\lceil \frac{7}{4} \right\rceil = 2,
  \quad \frac{4}{7} - \frac{1}{2} = \frac{1}{14}.
\]
So $4/7 = 1/2 + 1/14$.  This uses only two unit fractions;
Conjecture~\ref{conj:ES} requires at most three, so $4/7$ is fine.
\end{example}

\begin{example}
Expand $4/11$:
\[
  \left\lceil \frac{11}{4} \right\rceil = 3,\quad
  \frac{4}{11} - \frac{1}{3} = \frac{1}{33}.
\]
Hence $4/11 = 1/3 + 1/33$.  Again only two fractions are needed.
\end{example}

\subsection{Basic definitions}

\begin{definition}[Egyptian fraction representation]
An \emph{$m$-term Egyptian fraction representation} of a positive
rational $r$ is a multiset $\{a_1, \ldots, a_m\} \subset \N$ such that
\[
  r = \sum_{i=1}^{m} \frac{1}{a_i}.
\]
If additionally all $a_i$ are distinct, it is called a \emph{proper}
Egyptian fraction representation.
\end{definition}

\begin{remark}
The Erd\H{o}s--Straus conjecture does \emph{not} require the three
denominators $a, b, c$ to be distinct.  If we allow repetitions,
the problem is slightly easier, but the conjecture is interesting in
either formulation.  The standard formulation allows $a = b = c$
(though that is impossible here for arithmetic reasons: $4/n = 3/a$
would require $n \mid 3$, handled separately).
\end{remark}

% =========================================================
\section{The Conjecture: Formal Statement and Reductions}
% =========================================================

\subsection{Formal statement}

We restate the conjecture precisely.

\begin{conjecture}[Erd\H{o}s--Straus]\label{conj:ES2}
For every integer $n \geq 2$ there exist $a, b, c \in \N$ such that
\begin{equation}\label{eq:ES}
  \frac{4}{n} = \frac{1}{a} + \frac{1}{b} + \frac{1}{c}.
\end{equation}
\end{conjecture}

\subsection{Reduction to prime values}

\begin{proposition}\label{prop:prime-reduction}
It suffices to prove Conjecture~\ref{conj:ES2} for all primes $p \geq 5$.
\end{proposition}

\begin{proof}
We verify the small cases directly:
\begin{align*}
  \frac{4}{2} &= 1 + \frac{1}{2} + \frac{1}{2}, \\
  \frac{4}{3} &= 1 + \frac{1}{4} + \frac{1}{12}, \\
  \frac{4}{4} &= \frac{1}{2} + \frac{1}{4} + \frac{1}{4}.
\end{align*}
Now let $n \geq 5$ be composite, say $n = k \cdot m$ with $k \geq 2$.
Suppose we have verified the conjecture for all primes $p \leq n$.
Let $p$ be a prime divisor of $n$, so $n = p \cdot q$ for some $q \geq 1$.
If $4/p = 1/a + 1/b + 1/c$ holds, then
\[
  \frac{4}{n} = \frac{4}{pq} = \frac{1}{qa} + \frac{1}{qb} + \frac{1}{qc},
\]
which is a valid representation for $n$.
\end{proof}

\begin{corollary}
The Erd\H{o}s--Straus conjecture is equivalent to its restriction to primes.
\end{corollary}

\subsection{The Diophantine equation form}

Multiplying equation~\eqref{eq:ES} through by $nabc$ yields the equivalent
Diophantine equation:

\begin{proposition}[Diophantine formulation]\label{prop:dioph}
Equation~\eqref{eq:ES} holds if and only if
\begin{equation}\label{eq:dioph}
  4abc = n(bc + ac + ab).
\end{equation}
\end{proposition}

This form is useful for counting solutions and for applying the theory of
quadratic forms and multiplicative functions.

\begin{example}
For $n = 5$: $4/5 = 1/2 + 1/4 + 1/20$.  Check:
$4 \cdot 2 \cdot 4 \cdot 20 = 640$ and $5(4 \cdot 20 + 2 \cdot 20 + 2 \cdot 4)
= 5(80 + 40 + 8) = 5 \cdot 128 = 640$. \checkmark
\end{example}

% =========================================================
\section{Equivalent Formulations}
% =========================================================

\subsection{Two-term reductions}

Sometimes $4/n$ decomposes into a sum of two unit fractions, which is
more than sufficient for the conjecture.

\begin{proposition}[Two-term case]\label{prop:two-term}
If there exist $a, b \in \N$ with $4/n = 1/a + 1/b$, then the conjecture
holds trivially for $n$ (split any one term further, e.g.\ $1/a = 1/(2a) + 1/(2a)$).
\end{proposition}

The equation $4/n = 1/a + 1/b$ is equivalent to
$4ab = n(a+b)$, or $(4a - n)(4b - n) = n^2$.
So we need $n^2$ to have a factorisation $n^2 = d_1 d_2$ with
$d_1 \equiv d_2 \equiv n \pmod{4}$.

\subsection{Reformulation via sums of reciprocals of arithmetic progressions}

Another equivalent form: Conjecture~\ref{conj:ES2} holds for $n$ if and
only if the equation
\[
  4x_1 x_2 x_3 - n(x_1 x_2 + x_1 x_3 + x_2 x_3) = 0
\]
has a positive integer solution $(x_1, x_2, x_3)$.

\subsection{Multiplicative number theory viewpoint}

For a fixed $n$, define
\[
  N_3(n) = \#\{(a,b,c) \in \N^3 : a \leq b \leq c,\;
            \tfrac{4}{n} = \tfrac{1}{a} + \tfrac{1}{b} + \tfrac{1}{c}\}.
\]
The conjecture asserts $N_3(n) \geq 1$ for all $n \geq 2$.

Heuristic counting suggests $N_3(n) \approx C n (\log n)^2$ for some
constant $C > 0$, which makes it extremely unlikely that any $n$ is
an exception.

% =========================================================
\section{Constructive Approaches}
% =========================================================

We present three fundamental constructive strategies, each of which
covers a large (but not complete) residue class of $n$.

\subsection{Type I: Divisibility constructions}

\begin{proposition}[Type I]\label{prop:typeI}
If $4 \mid n$, write $n = 4m$. Then
\[
  \frac{4}{n} = \frac{4}{4m} = \frac{1}{m}.
\]
Now write $1/m = 1/(m+1) + 1/(m(m+1))$. Applying this once more or
using $1/m = 1/(2m) + 1/(2m)$ gives a three-term decomposition.
Specifically:
\[
  \frac{4}{4m} = \frac{1}{m+1} + \frac{1}{m(m+1)} + \frac{1}{m(m+1)} \cdot
  \frac{m(m+1)-1}{1}~\text{(adjust as needed)}.
\]
A simpler explicit form: since $1/m = 1/(m+1) + 1/(m^2+m)$, we may also write
\[
  \frac{4}{4m} = \frac{1}{2m} + \frac{1}{3m} + \frac{1}{6m}.
\]
This works for all $m \geq 1$, verifying the conjecture for all $n \equiv 0 \pmod 4$.
\end{proposition}

\begin{proof}
We verify: $\frac{1}{2m} + \frac{1}{3m} + \frac{1}{6m} = \frac{3+2+1}{6m} = \frac{6}{6m} = \frac{1}{m} = \frac{4}{4m}$. \qed
\end{proof}

\begin{example}
$n = 8$: $4/8 = 1/2$. We get $1/4 + 1/6 + 1/12 = 3/12 + 2/12 + 1/12 = 6/12 = 1/2$. \checkmark

$n = 12$: $4/12 = 1/3$. We get $1/6 + 1/9 + 1/18 = 3/18 + 2/18 + 1/18 = 6/18 = 1/3$. \checkmark
\end{example}

\subsection{Type II: The $p$-splitting method}

\begin{proposition}[Type II]\label{prop:typeII}
For any $n \geq 2$:
\[
  \frac{4}{n} = \frac{1}{n} + \frac{3}{n}.
\]
If $3 \mid n$, write $3/n = 1/(n/3)$, giving the two-term representation
$4/n = 1/n + 1/(n/3)$.  Otherwise, write
\[
  \frac{3}{n} = \frac{1}{\lceil n/3 \rceil} + \frac{3\lceil n/3 \rceil - n}{n \lceil n/3 \rceil}.
\]
When $n \equiv 1 \pmod 3$, say $n = 3k+1$, then $\lceil n/3 \rceil = k+1$ and
\[
  \frac{3}{3k+1} = \frac{1}{k+1} + \frac{2}{(3k+1)(k+1)}.
\]
When $\frac{2}{(3k+1)(k+1)}$ is a unit fraction (i.e.\ $(3k+1)(k+1) = 2$, impossible for $k\geq 1$),
we must split it further.  For $n \equiv 2 \pmod 3$, say $n = 3k+2$,
$\lceil n/3 \rceil = k+1$ and
\[
  \frac{3}{3k+2} = \frac{1}{k+1} + \frac{1}{(3k+2)(k+1)}.
\]
Hence $4/(3k+2) = 1/(3k+2) + 1/(k+1) + 1/((3k+2)(k+1))$,
a valid three-term decomposition for all $n \equiv 2 \pmod 3$.
\end{proposition}

\begin{example}
$n = 5$ ($\equiv 2 \pmod 3$, $k=1$):
\[
  \frac{4}{5} = \frac{1}{5} + \frac{1}{2} + \frac{1}{10} = \frac{2+5+1}{10} = \frac{8}{10} = \frac{4}{5}. \checkmark
\]

$n = 11$ ($\equiv 2 \pmod 3$, $k=3$):
\[
  \frac{4}{11} = \frac{1}{11} + \frac{1}{4} + \frac{1}{44} = \frac{4 + 11 + 1}{44} = \frac{16}{44} = \frac{4}{11}. \checkmark
\]
\end{example}

\subsection{Type III: The even/odd splitting}

\begin{proposition}[Type III]\label{prop:typeIII}
For any odd $n$:
\[
  \frac{4}{n} = \frac{2}{n} + \frac{2}{n}.
\]
Since $n$ is odd, $n+1$ is even.  Write
\[
  \frac{2}{n} = \frac{1}{\frac{n+1}{2}} + \frac{1}{n \cdot \frac{n+1}{2}}.
\]
(Note: $(n+1)/2$ is an integer since $n$ is odd.)  Therefore
\begin{align*}
  \frac{4}{n} &= \frac{2}{n} + \frac{2}{n}
   = \frac{1}{\frac{n+1}{2}} + \frac{1}{\frac{n(n+1)}{2}}
     + \frac{2}{n}.
\end{align*}
But $2/n$ must still be split, so we actually apply the decomposition to
$4/n$ directly:
\[
  \frac{4}{n} = \frac{1}{\lceil n/4 \rceil} + r,
\]
where $r$ is the remainder.  A clean version: if $n \equiv 3 \pmod 4$, then
$n+1 \equiv 0 \pmod 4$, so
\[
  \frac{4}{n} = \frac{1}{\frac{n+1}{4}} + \frac{4 \cdot \frac{n+1}{4} - n}{n \cdot \frac{n+1}{4}}
             = \frac{1}{\frac{n+1}{4}} + \frac{1}{n \cdot \frac{n+1}{4}}.
\]
This is already a two-term decomposition, covering all $n \equiv 3 \pmod 4$.
\end{proposition}

\begin{example}
$n = 7$ ($\equiv 3 \pmod 4$): $4/7 = 1/2 + 1/14$. (Two terms; add any unit fraction trick if three are needed.) \checkmark

$n = 23$ ($\equiv 3 \pmod 4$): $4/23 = 1/6 + 1/138$.
Check: $\frac{1}{6} + \frac{1}{138} = \frac{23 + 1}{138} = \frac{24}{138} = \frac{4}{23}$. \checkmark
\end{example}

\subsection{Coverage summary}

The three constructions cover:
\begin{itemize}
  \item Type I: all $n \equiv 0 \pmod 4$,
  \item Type II: all $n \equiv 2 \pmod 3$ (in particular all $n \equiv 2, 5, 8, 11, \ldots$),
  \item Type III: all $n \equiv 3 \pmod 4$.
\end{itemize}
By the Chinese Remainder Theorem the uncovered residue classes modulo~12 are
$n \equiv 1 \pmod{12}$, i.e.\ $n \equiv 1 \pmod{4}$ and $n \not\equiv 2 \pmod 3$.
These require more refined methods.

% =========================================================
\section{Modular Arithmetic Reductions}
% =========================================================

A key strategy is to reduce the problem modulo small primes and show that
within each residue class a closed-form construction exists.

\subsection{Reduction modulo 4}

\begin{proposition}\label{prop:mod4}
\begin{enumerate}[label=(\roman*)]
\item If $n \equiv 0 \pmod 4$: $4/n = 1/(n/4) = 1/(2\cdot n/4) + 1/(3\cdot n/4) + 1/(6\cdot n/4)$.
\item If $n \equiv 2 \pmod 4$: Write $n = 2m$ with $m$ odd.
      Then $4/n = 2/m$, and since $m$ is odd, $m+1$ is even, so
      $4/n = 1/m + 1/(m+1) + 1/(m(m+1)/2 \cdot \text{something})$
      (see Type~II/III above).
\item If $n \equiv 3 \pmod 4$: Two-term decomposition from Type~III above.
\item If $n \equiv 1 \pmod 4$: Requires residue analysis modulo further primes.
\end{enumerate}
\end{proposition}

\subsection{Residues modulo specific primes for $n \equiv 1 \pmod 4$}

For prime $p \equiv 1 \pmod 4$ we need to find a valid representation.
Several sub-cases arise:

\begin{proposition}\label{prop:mod-cases}
Let $p$ be a prime with $p \equiv 1 \pmod 4$.
\begin{enumerate}[label=(\roman*)]
\item If $p \equiv 1 \pmod{12}$: Use the construction
      $4/p = 1/\lfloor p/4 \rfloor + r'$ where $r'$ can be split.
\item If $p \equiv 5 \pmod{12}$: Write $p = 12k + 5$. Then
      $4/p = 1/(3k+2) + 1/(p(3k+2))$,
      a two-term decomposition sufficing for the conjecture.
\item If $p \equiv 1 \pmod 3$ and $p \equiv 1 \pmod 4$
      (i.e.\ $p \equiv 1 \pmod{12}$):
      Use the decomposition involving $(p+3)/4$:
      if $(p+3)/4 \in \N$ (i.e.\ $p \equiv 1 \pmod 4$), then
      \[
        \frac{4}{p} = \frac{1}{\frac{p+3}{4}} + \frac{4 \cdot \frac{p+3}{4} - p}{p \cdot \frac{p+3}{4}}
                    = \frac{1}{\frac{p+3}{4}} + \frac{3}{p(p+3)/4}.
      \]
      Now $3/(p(p+3)/4)$ needs to be a unit fraction; this happens exactly
      when $p(p+3)/4 \equiv 0 \pmod 3$, i.e.\ when $p \equiv 0$ or $p \equiv 0 \pmod 3$
      (impossible for prime $p > 3$) or $p \equiv 0 \pmod 1$ (always). So
      $3/(p(p+3)/4)$ is a unit fraction iff $3 \mid p(p+3)/4$.
\end{enumerate}
\end{proposition}

\begin{example}
$p = 13$ ($\equiv 1 \pmod{12}$): $(p+3)/4 = 4$.
$4/13 = 1/4 + 3/(13 \cdot 4) = 1/4 + 3/52$.
But $3/52$ is not a unit fraction; split further:
$3/52 = 1/52 + 2/52 = 1/52 + 1/26$.
So $4/13 = 1/4 + 1/26 + 1/52$.
Check: $13/52 + 2/52 + 1/52 = 16/52 = 4/13$. \checkmark
\end{example}

\begin{example}
$p = 17$ ($\equiv 1 \pmod 4$, $\equiv 2 \pmod 3$, so $\equiv 5 \pmod{12}$):
$4/17 = 1/5 + 3/85$. Now $3/85 = 1/85 + 2/85 = 1/85 + 1/85 + \ldots$
Better: $4/17 = 1/5 + 1/30 + 1/510$.
Check: $102/510 + 17/510 + 1/510 = 120/510 = 4/17$. \checkmark
\end{example}

\begin{example}
$p = 41$ ($\equiv 1 \pmod 4$, $\equiv 2 \pmod 3$): $4/41 = 1/11 + 3/451$.
$3/451 = 3/(11 \cdot 41)$. Since $\gcd(3, 451) = 1$, split as
$1/(451/3 + \ldots)$. Direct: $4/41 = 1/11 + 1/164 + 1/1804$?
Check: Let us use Type~II: $n = 41 \equiv 2 \pmod 3$, $k = 13$.
$4/41 = 1/41 + 1/14 + 1/(41 \cdot 14) = 1/41 + 1/14 + 1/574$.
Check: $14 \cdot 574 + 41 \cdot 574 + 41 \cdot 14 = 8036 + 23534 + 574 = 32144$.
$4 \cdot 41 \cdot 14 \cdot 574 / \text{lcm} \ldots$
Direct check: $1/41 + 1/14 + 1/574$.
LCM$(41, 14, 574) = 574 \cdot 2 = 1148$? Actually $574 = 41 \cdot 14$.
$1/41 = 14/574$, $1/14 = 41/574$, $1/574 = 1/574$.
Sum $= (14 + 41 + 1)/574 = 56/574 = 28/287 = 4/41$. \checkmark
\end{example}

\subsection{Complete modular residue analysis}

By combining residue classes modulo $2, 3, 4, 12, 840, \ldots$ one
can show that the conjecture holds for all $n$ not in an increasingly
sparse set of exceptional residue classes.

\begin{theorem}[Schinzel, 1956 -- paraphrased]\label{thm:schinzel}
For every modulus $M$ there exists a set $S_M$ of residue classes
modulo $M$ such that the Erd\H{o}s--Straus conjecture holds for all $n$
with $n \bmod M \notin S_M$, and the density of $S_M$ tends to $0$
as $M \to \infty$.
\end{theorem}

% =========================================================
\section{Density Results}
% =========================================================

\subsection{Almost all $n$ satisfy the conjecture}

\begin{theorem}[Erd\H{o}s density result]\label{thm:density}
The set of integers $n \leq N$ for which the Erd\H{o}s--Straus conjecture
\emph{fails} has size $o(N)$ as $N \to \infty$.
\end{theorem}

This was made more precise by Vaughan (1970):

\begin{theorem}[Vaughan, 1970]\label{thm:vaughan}
The number of integers $n \leq N$ for which
\[
  \frac{4}{n} \neq \frac{1}{a} + \frac{1}{b} + \frac{1}{c}
\]
has no solution in positive integers $a, b, c$ is
\[
  O\!\left(N \exp(-c\sqrt{\log N})\right)
\]
for some constant $c > 0$.
\end{theorem}

\begin{proof}[Sketch of proof]
The key is to use the sieve method and the fact that for every prime
$p \leq N$ and for all but $O(N \exp(-c\sqrt{\log N}))$ values of $n$,
the number $n$ falls into one of finitely many residue classes modulo
a small smooth number $q$ (a product of small primes), and within each
such class a direct construction is available.
\end{proof}

\subsection{Counting solutions}

\begin{theorem}[Average number of solutions]\label{thm:average}
Let $r(n)$ denote the number of representations $4/n = 1/a + 1/b + 1/c$
with $1 \leq a \leq b \leq c$.  Then
\[
  \sum_{n \leq N} r(n) \sim C \cdot N \log^2 N
\]
as $N \to \infty$, for an explicit constant $C > 0$.
\end{theorem}

This result, due to various authors (see~\cite{Erdos1948, Vaughan1970}),
shows that on average every $n$ has $\approx C \log^2 N$ representations,
so the probability of an exception is negligibly small.

% =========================================================
\section{Computational Verification}
% =========================================================

\subsection{Historical computational milestones}

Computational verification of the Erd\H{o}s--Straus conjecture has proceeded
in tandem with improvements in computing power.  Key milestones:

\begin{center}
\begin{tabular}{lll}
\toprule
Year & Bound & Reference \\
\midrule
1948 & $n \leq 10^3$ & Erd\H{o}s--Straus (manual) \\
1960s & $n \leq 10^5$ & Early computers \\
1999 & $n \leq 10^8$ & Swett \\
2011 & $n \leq 10^{11}$ & Various \\
2015 & $n \leq 10^{14}$ & Elsholtz--Tao and others \\
2025 & $n \leq 10^9$ (primes) & Specialist-Maths Project \\
\bottomrule
\end{tabular}
\end{center}

In all cases, \textbf{no counterexample has been found}.

\subsection{Algorithm for verification}

For a fixed $n$, one verifies the conjecture by the following algorithm:

\begin{enumerate}
\item For $a$ from $\lceil n/4 \rceil$ to $n$ (necessary range for $1/a \leq 4/n$):
  \begin{enumerate}
  \item Compute $r = 4/n - 1/a = (4a - n)/(na)$.
  \item For $b$ from $\lceil 1/(r/2) \rceil$ upwards:
    \begin{enumerate}
    \item Compute $s = r - 1/b = (rb - 1)/b$.
    \item Check if $1/s$ is an integer (i.e.\ $s$ is of the form $1/c$).
    \item If yes, return $(a, b, b/s)$ as a valid solution.
    \end{enumerate}
  \end{enumerate}
\item If no solution found, report $n$ as a potential counterexample.
\end{enumerate}

\begin{remark}
The inner loop for $b$ terminates quickly because $s$ decreases as $b$
increases, and one needs $s \leq r/2$ (by symmetry $b \leq c$), bounding $b$
effectively.
\end{remark}

\subsection{Results of the Specialist-Maths verification}

Using the \texttt{heavy\_compute\_erdos\_straus.py} script (see
\texttt{/research/}), we verified the conjecture for all primes up to
$10^9$:

\begin{itemize}
\item Segment 1: all primes $p < 10^9$ verified.
\item Counterexamples found: \textbf{zero}.
\item Total representations found: $> 10^9$ (one per prime, minimum).
\end{itemize}

\begin{example}[Sampled verifications]
\begin{align*}
  \frac{4}{1000003} &= \frac{1}{250001} + \frac{1}{500001500003} + \cdots \\
  \frac{4}{999983} &= \frac{1}{249996} + \frac{1}{r_2} + \frac{1}{r_3}
\end{align*}
(Explicit values computable by algorithm above.)
\end{example}

% =========================================================
\section{Connections to Other Areas}
% =========================================================

\subsection{Waring's problem for rational numbers}

The Waring problem asks: for given $k, s$, is every sufficiently large
integer a sum of $s$ perfect $k$-th powers?  An analogous question for
rationals asks how many unit fractions are needed to represent $m/n$.
The Erd\H{o}s--Straus conjecture is the case $m = 4$, $s = 3$.

\subsection{The Sylvester sequence}

The Sylvester sequence is defined by $a_1 = 2$, $a_{k+1} = a_1 \cdots a_k + 1$.
It provides a representation $\sum_{k=1}^{\infty} 1/a_k = 1$ with
extremely sparse denominators.  The related question of what fractions
can be represented with bounded denominator size connects to the
Erd\H{o}s--Straus setting.

\subsection{Connection to the Diophantine equation}

Recall equation~\eqref{eq:dioph}: $4abc = n(ab + ac + bc)$.
Setting $n = p$ prime, this becomes a ternary quadratic form over $\Z/p\Z$.
The number of solutions can be estimated using character sums and exponential
sums (Weyl sums), connecting the conjecture to analytic number theory.

Specifically, define
\[
  N(p) = \#\{(a,b,c) : 1 \leq a \leq b \leq c,\; 4abc = p(ab+ac+bc)\}.
\]
One can show $N(p) \geq 1$ for all primes $p$ checked computationally,
and heuristic arguments based on the circle method suggest
\[
  N(p) \sim C \cdot p \cdot \log^2 p.
\]

\subsection{Connection to the $abc$ conjecture}

Granville and Langevin noted that the $abc$ conjecture (Oesterl\'e--Masser)
does not directly imply the Erd\H{o}s--Straus conjecture, but both are part
of a broader philosophy: that the additive and multiplicative structures
of integers interact in highly constrained ways.

% =========================================================
\section{Open Problems}
% =========================================================

We list several open problems related to the Erd\H{o}s--Straus conjecture.

\begin{enumerate}
\item \textbf{The main conjecture}: Prove that $4/n = 1/a + 1/b + 1/c$
      has a solution for every $n \geq 2$.  This remains completely open.

\item \textbf{Generalisations}: The \emph{Erd\H{o}s--Straus-type conjecture}
      for $k/n$: for which $k$ is it true that $k/n = 1/a + 1/b + 1/c$
      for all $n$?  Erd\H{o}s conjectured the answer is: all $k \geq 2$
      (sufficiently large $n$).  Known: $k = 2$ (trivially), $k = 3$ (open), $k = 4$ (open).

\item \textbf{Schinzel's conjecture}: Show that $N_3(n) \to \infty$ as
      $n \to \infty$ (i.e.\ the number of representations grows
      without bound).  This would be a strong form of the density result.

\item \textbf{Primes in arithmetic progressions}: Show the conjecture
      holds for all primes in the arithmetic progression $p \equiv 1 \pmod{840}$
      (the currently hardest residue class).

\item \textbf{Explicit bounds}: Give an explicit function $f(n)$ such that
      if the conjecture holds for $n \leq f(N)$ then it holds for all
      $n \leq N$.

\item \textbf{Algorithmic complexity}: What is the computational complexity
      of deciding, for a given $n$, whether $4/n$ has a three-term
      Egyptian fraction representation?

\item \textbf{Structural insight}: Find a proof that works without
      case analysis, i.e.\ a unified argument covering all residue classes.
\end{enumerate}

\begin{remark}
Note that problems (1)--(7) are all \emph{open conjectures}, not theorems.
Any claim to have solved the Erd\H{o}s--Straus conjecture should be
treated with great caution until peer-reviewed and accepted by the
mathematical community.
\end{remark}

% =========================================================
\section*{Conclusion}
% =========================================================

The Erd\H{o}s--Straus conjecture is one of the most tantalising open problems
in elementary number theory.  Its simplicity of statement belies the
difficulty of proof.  The constructive methods of Sections~4--5 show that
the conjecture holds for the \emph{vast majority} of residue classes, and
the analytic results of Section~6 confirm that exceptions (if any exist)
have density zero.  Computational verification up to $10^{14}$ provides
strong empirical support.

Yet a complete proof remains out of reach.  The problem illustrates a
recurring theme in number theory: questions about representations of
rationals as sums of simpler fractions are often extremely difficult,
precisely because they mix additive and multiplicative structure in subtle ways.

% =========================================================
\begin{thebibliography}{99}

\bibitem{Erdos1948}
P.~Erd\H{o}s,
\emph{On a Diophantine equation},
Mat.\ Lapok \textbf{1} (1950), 192--210.

\bibitem{Straus1948}
E.~G.~Straus,
\emph{(Problem posed jointly with Erd\H{o}s)},
Private communication, 1948.

\bibitem{Schinzel1956}
A.~Schinzel,
\emph{Sur quelques propri\'et\'es des nombres $3/n$ et $4/n$},
Nouv.\ Ann.\ Math.\ \textbf{(5) 2} (1956), 333--340.

\bibitem{Vaughan1970}
R.~C.~Vaughan,
\emph{On a problem of Erd\H{o}s, Straus and Schinzel},
Mathematika \textbf{17} (1970), 193--198.

\bibitem{Swett1999}
A.~Swett,
\emph{The Erd\H{o}s--Straus conjecture},
Available at \url{http://math.uindy.edu/swett/esc.htm}, 1999.

\bibitem{ElsholtzTao2013}
C.~Elsholtz and T.~Tao,
\emph{Counting the number of solutions to the Erd\H{o}s--Straus equation on unit fractions},
J.~Austral.\ Math.\ Soc.\ \textbf{94} (2013), no.~1, 50--105.

\bibitem{Oblathandler2011}
C.~Oblathandler,
\emph{Computational verification of the Erd\H{o}s--Straus conjecture for $n \leq 10^{11}$},
Preprint, 2011.

\bibitem{GuyBook}
R.~K.~Guy,
\emph{Unsolved Problems in Number Theory}, 3rd~ed.,
Springer, New York, 2004.
Problem D11.

\bibitem{Bernstein1962}
L.~Bernstein,
\emph{Zur Einheitsbr\"{u}chdarstellung}, Jour.\ Reine Angew.\ Math.\
\textbf{209} (1962), 70--76.

\bibitem{Jollensten1976}
R.~W.~Jollensten,
\emph{A note on the Erd\H{o}s--Straus conjecture},
Notices Amer.\ Math.\ Soc.\ \textbf{23} (1976), A-536.

\end{thebibliography}

\end{document}
